% GENERATED FILE. Edit canonical lessons, then run npm run generate:books.
% canonical-source-hash: fnv1a64:927d8a780a01ce60
% canonical-lessons: SA-C363-adhyayanam, SA-C363-visayah, SA-C363-sastram, SA-C363-adhite, SA-C363-abhyasyati, SA-R363-review-356-363

\chapter{Studying a Subject}
\label{ch:sa-studying-a-subject}
\hlchaptermodality{363}

\begin{chapteropening}
\textbf{By the end of this chapter:} \emph{I can say what I study and how I practise it.}

Complete the last lesson of chapter 363: I can say what I study and how I practise it.
\end{chapteropening}

\section[\texorpdfstring{\sk{अध्ययनम्} (adhyayanam)}{अध्ययनम् (adhyayanam)}]{\sk{अध्ययनम्} (adhyayanam) — studying; study}

\label{lesson:SA-C363-adhyayanam}



\begin{quote}
Before the new one: say the Sanskrit for inspects, then the Sanskrit for devises.
\end{quote}

\subsection*{You'll want to know: \sk{अध्ययनम्}}
\textbf{\sk{अध्ययनम्}} — \emph{adhyayanam} — \textquotedblleft{}studying; study\textquotedblright{}.

\begin{quote}
\sk{मम} \sk{अध्ययनं} \sk{प्रातः} \sk{भवति।}
\end{quote}

\begin{quote}
\emph{mama adhyayanaṃ prātaḥ bhavati.}
\end{quote}

\begin{quote}
I do my studying in the morning.
\end{quote}

A neuter noun.

\begin{cousinweb}
\textbf{\sk{अध्ययनम्}} (\emph{adhyayanam}) is built on \emph{adhi}, \textquotedblleft{}over, upon\textquotedblright{}, and the root √i, \textquotedblleft{}to go\textquotedblright{}; MW gives \emph{adhy-ayana} \textquotedblleft{}reading, studying\textquotedblright{} (MW).
\end{cousinweb}

\begin{grammarlens}[title={What you've built}]
One more word for giving reasons and opinions.
\end{grammarlens}

\subsection*{Your turn}
\begin{itemize}
\raggedright
  \item \emph{Say it:} \emph{adhyayanam}
  \item \emph{Say it:} \emph{adhyayanam}, once more
  \item \emph{Say it:} \emph{mama adhyayanaṃ prātaḥ bhavati.} — I do my studying in the morning.
  \item \emph{Recall:} say the Sanskrit for inspects, then the Sanskrit for devises, then say \emph{adhyayanam} again
\end{itemize}

\begin{tcolorbox}[breakable,colback=teal!4,colframe=teal!35!black,title={Before you move on}]
What does \sk{अध्ययनम्} mean? (\textquotedblleft{}Studying; study\textquotedblright{}.) Say it once more.
\end{tcolorbox}
\section[\texorpdfstring{\sk{विषयः} (viṣayaḥ)}{विषयः (viṣayaḥ)}]{\sk{विषयः} (viṣayaḥ) — a topic; a subject (of study)}

\label{lesson:SA-C363-visayah}



\begin{quote}
Before the new one: say the Sanskrit for devises, then the Sanskrit for studying.
\end{quote}

\subsection*{You'll want to know: \sk{विषयः}}
\textbf{\sk{विषयः}} — \emph{viṣayaḥ} — \textquotedblleft{}a topic; a subject (of study)\textquotedblright{}.

\begin{quote}
\sk{मम} \sk{प्रियः} \sk{विषयः} \sk{इतिहासः} \sk{अस्ति।}
\end{quote}

\begin{quote}
\emph{mama priyaḥ viṣayaḥ itihāsaḥ asti.}
\end{quote}

\begin{quote}
My favourite subject is history.
\end{quote}

A masculine noun.

\begin{cousinweb}
\textbf{\sk{विषयः}} (\emph{viṣayaḥ}) MW gives \emph{viṣaya} \textquotedblleft{}an object of sense\textquotedblright{} and \textquotedblleft{}any object of concern or attention\textquotedblright{}, and \textquotedblleft{}a subject of study\textquotedblright{} is how modern Sanskrit uses it; MW gives no sure derivation, so its origin is uncertain (MW).
\end{cousinweb}

\begin{grammarlens}[title={What you've built}]
One more word for giving reasons and opinions.
\end{grammarlens}

\subsection*{Your turn}
\begin{itemize}
\raggedright
  \item \emph{Say it:} \emph{viṣayaḥ}
  \item \emph{Say it:} \emph{viṣayaḥ}, once more
  \item \emph{Say it:} \emph{mama priyaḥ viṣayaḥ itihāsaḥ asti.} — My favourite subject is history.
  \item \emph{Recall:} say the Sanskrit for devises, then the Sanskrit for studying, then say \emph{viṣayaḥ} again
\end{itemize}

\begin{tcolorbox}[breakable,colback=teal!4,colframe=teal!35!black,title={Before you move on}]
What does \sk{विषयः} mean? (\textquotedblleft{}A topic; a subject (of study)\textquotedblright{}.) Say it once more.
\end{tcolorbox}
\section[\texorpdfstring{\sk{शास्त्रम्} (śāstram)}{शास्त्रम् (śāstram)}]{\sk{शास्त्रम्} (śāstram) — a branch of learning; a treatise}

\label{lesson:SA-C363-sastram}



\begin{quote}
Before the new one: say the Sanskrit for studying, then the Sanskrit for a topic.
\end{quote}

\subsection*{You'll want to know: \sk{शास्त्रम्}}
\textbf{\sk{शास्त्रम्}} — \emph{śāstram} — \textquotedblleft{}a branch of learning; a treatise\textquotedblright{}.

\begin{quote}
\sk{पण्डितः} \sk{शास्त्रं} \sk{पाठयति।}
\end{quote}

\begin{quote}
\emph{paṇḍitaḥ śāstraṃ pāṭhayati.}
\end{quote}

\begin{quote}
The scholar teaches a branch of learning.
\end{quote}

A neuter noun.

\begin{cousinweb}
\textbf{\sk{शास्त्रम्}} (\emph{śāstram}) comes from the root √śās, \textquotedblleft{}to instruct, to order\textquotedblright{}; MW gives \emph{śāstra} \textquotedblleft{}teaching, instruction\textquotedblright{} and \textquotedblleft{}any book or treatise\textquotedblright{} (MW).
\end{cousinweb}

\begin{grammarlens}[title={What you've built}]
One more word for giving reasons and opinions.
\end{grammarlens}

\subsection*{Your turn}
\begin{itemize}
\raggedright
  \item \emph{Say it:} \emph{śāstram}
  \item \emph{Say it:} \emph{śāstram}, once more
  \item \emph{Say it:} \emph{paṇḍitaḥ śāstraṃ pāṭhayati.} — The scholar teaches a branch of learning.
  \item \emph{Recall:} say the Sanskrit for studying, then the Sanskrit for a topic, then say \emph{śāstram} again
\end{itemize}

\begin{tcolorbox}[breakable,colback=teal!4,colframe=teal!35!black,title={Before you move on}]
What does \sk{शास्त्रम्} mean? (\textquotedblleft{}A branch of learning; a treatise\textquotedblright{}.) Say it once more.
\end{tcolorbox}
\section[\texorpdfstring{\sk{अधीते} (adhīte)}{अधीते (adhīte)}]{\sk{अधीते} (adhīte) — studies}

\label{lesson:SA-C363-adhite}



\begin{quote}
Before the new one: say the Sanskrit for a topic, then the Sanskrit for a branch of learning.
\end{quote}

\subsection*{You'll want to know: \sk{अधीते}}
\textbf{\sk{अधीते}} — \emph{adhīte} — \textquotedblleft{}studies\textquotedblright{}.

\begin{quote}
\sk{अरुणः} \sk{विश्वविद्यालये} \sk{विज्ञानम्} \sk{अधीते।}
\end{quote}

\begin{quote}
\emph{aruṇaḥ viśvavidyālaye vijñānam adhīte.}
\end{quote}

\begin{quote}
Arun studies science at the university.
\end{quote}

Past (he …): \sk{सः} \sk{गणितम्} \sk{अधीतवान्}, \emph{saḥ gaṇitam adhītavān}.

\begin{cousinweb}
\textbf{\sk{अधीते}} (\emph{adhīte}) is built on \emph{adhi}, \textquotedblleft{}over\textquotedblright{}, and the root √i, \textquotedblleft{}to go\textquotedblright{}; MW gives \emph{adhī} \textquotedblleft{}to study, to learn by heart, to read\textquotedblright{} (MW).
\end{cousinweb}

\begin{grammarlens}[title={What you've built}]
One more word for giving reasons and opinions.
\end{grammarlens}

\subsection*{Your turn}
\begin{itemize}
\raggedright
  \item \emph{Say it:} \emph{adhīte}
  \item \emph{Say it:} \emph{adhīte}, once more
  \item \emph{Say it:} \emph{aruṇaḥ viśvavidyālaye vijñānam adhīte.} — Arun studies science at the university.
  \item \emph{Recall:} say the Sanskrit for a topic, then the Sanskrit for a branch of learning, then say \emph{adhīte} again
\end{itemize}

\begin{tcolorbox}[breakable,colback=teal!4,colframe=teal!35!black,title={Before you move on}]
What does \sk{अधीते} mean? (\textquotedblleft{}Studies\textquotedblright{}.) Say it once more.
\end{tcolorbox}
\section[\texorpdfstring{\sk{अभ्यस्यति} (abhyasyati)}{अभ्यस्यति (abhyasyati)}]{\sk{अभ्यस्यति} (abhyasyati) — practises (a skill)}

\label{lesson:SA-C363-abhyasyati}



\begin{quote}
Before the new one: say the Sanskrit for a branch of learning, then the Sanskrit for studies.
\end{quote}

\subsection*{You'll want to know: \sk{अभ्यस्यति}}
\textbf{\sk{अभ्यस्यति}} — \emph{abhyasyati} — \textquotedblleft{}practises (a skill)\textquotedblright{}.

\begin{quote}
\sk{मीरा} \sk{प्रतिदिनं} \sk{वीणाम्} \sk{अभ्यस्यति।}
\end{quote}

\begin{quote}
\emph{mīrā pratidinaṃ vīṇām abhyasyati.}
\end{quote}

\begin{quote}
Mira practises the veena every day.
\end{quote}

Past (she …): \sk{सा} \sk{वीणाम्} \sk{अभ्यस्तवती}, \emph{sā vīṇām abhyastavatī}.

\begin{cousinweb}
\textbf{\sk{अभ्यस्यति}} (\emph{abhyasyati}) is built on \emph{abhi}, \textquotedblleft{}towards\textquotedblright{}, and the root √as, \textquotedblleft{}to throw\textquotedblright{}; MW gives \emph{abhy-√as} \textquotedblleft{}to practise, to exercise, to study\textquotedblright{} (MW).
\end{cousinweb}

\begin{grammarlens}[title={What you've built}]
One more word for giving reasons and opinions.
\end{grammarlens}

\subsection*{Your turn}
\begin{itemize}
\raggedright
  \item \emph{Say it:} \emph{abhyasyati}
  \item \emph{Say it:} \emph{abhyasyati}, once more
  \item \emph{Say it:} \emph{mīrā pratidinaṃ vīṇām abhyasyati.} — Mira practises the veena every day.
  \item \emph{Recall:} say the Sanskrit for a branch of learning, then the Sanskrit for studies, then say \emph{abhyasyati} again
\end{itemize}

\begin{tcolorbox}[breakable,colback=teal!4,colframe=teal!35!black,title={Before you move on}]
What does \sk{अभ्यस्यति} mean? (\textquotedblleft{}Practises (a skill)\textquotedblright{}.) Say it once more.
\end{tcolorbox}
\section[(dialogue)]{Looking back: from Notices from an Office to Studying a Subject}

\label{lesson:SA-R363-review-356-363}



\begin{quote}
Say the Sanskrit for notifies and appoints.
\end{quote}

\subsection*{Your turn}
Say these aloud:

\begin{itemize}
\raggedright
  \item notifies, appoints, approves of, a department, a proposal
  \item a committee, a council, foremost, a secretary, directs
  \item conduct, character, a good person, the right path, holds in regard
  \item renunciation, helping others, self-interest, craving, dedicates
  \item glory, standing, arrogance, an insult, misleads
  \item faith, steadfastness, hypocrisy, ungrateful, gives one's word
  \item research, an experiment, an inquiry, inspects, devises
  \item studying, a topic, a branch of learning, studies, practises
\end{itemize}

\begin{tcolorbox}[breakable,colback=teal!4,colframe=teal!35!black,title={Before you move on}]
Notifies? (\textbf{\sk{विज्ञापयति}}, \emph{vijñāpayati}.) Appoints? (\textbf{\sk{नियोजयति}}, \emph{niyojayati}.) Approves of? (\textbf{\sk{अनुमोदते}}, \emph{anumodate}.) A department? (\textbf{\sk{विभागः}}, \emph{vibhāgaḥ}.) A proposal? (\textbf{\sk{प्रस्तावः}}, \emph{prastāvaḥ}.) A committee? (\textbf{\sk{समितिः}}, \emph{samitiḥ}.) A council? (\textbf{\sk{परिषद्}}, \emph{pariṣad}.) Foremost? (\textbf{\sk{प्रमुखः}}, \emph{pramukhaḥ}.) A secretary? (\textbf{\sk{सचिवः}}, \emph{sacivaḥ}.) Directs? (\textbf{\sk{निर्दिशति}}, \emph{nirdiśati}.) Conduct? (\textbf{\sk{आचारः}}, \emph{ācāraḥ}.) Character? (\textbf{\sk{शीलम्}}, \emph{śīlam}.) A good person? (\textbf{\sk{सज्जनः}}, \emph{sajjanaḥ}.) The right path? (\textbf{\sk{सन्मार्गः}}, \emph{sanmārgaḥ}.) Holds in regard? (\textbf{\sk{आद्रियते}}, \emph{ādriyate}.) Renunciation? (\textbf{\sk{त्यागः}}, \emph{tyāgaḥ}.) Helping others? (\textbf{\sk{परोपकारः}}, \emph{paropakāraḥ}.) Self-interest? (\textbf{\sk{स्वार्थः}}, \emph{svārthaḥ}.) Craving? (\textbf{\sk{तृष्णा}}, \emph{tṛṣṇā}.) Dedicates? (\textbf{\sk{समर्पयति}}, \emph{samarpayati}.) Glory? (\textbf{\sk{यशः}}, \emph{yaśaḥ}.) Standing? (\textbf{\sk{प्रतिष्ठा}}, \emph{pratiṣṭhā}.) Arrogance? (\textbf{\sk{अहंकारः}}, \emph{ahaṅkāraḥ}.) An insult? (\textbf{\sk{अपमानः}}, \emph{apamānaḥ}.) Misleads? (\textbf{\sk{प्रतारयति}}, \emph{pratārayati}.) Faith? (\textbf{\sk{श्रद्धा}}, \emph{śraddhā}.) Steadfastness? (\textbf{\sk{निष्ठा}}, \emph{niṣṭhā}.) Hypocrisy? (\textbf{\sk{दम्भः}}, \emph{dambhaḥ}.) Ungrateful? (\textbf{\sk{कृतघ्नः}}, \emph{kṛtaghnaḥ}.) Gives one's word? (\textbf{\sk{प्रतिशृणोति}}, \emph{pratiśṛṇoti}.) Research? (\textbf{\sk{अनुसन्धानम्}}, \emph{anusandhānam}.) An experiment? (\textbf{\sk{प्रयोगः}}, \emph{prayogaḥ}.) An inquiry? (\textbf{\sk{अन्वेषणम्}}, \emph{anveṣaṇam}.) Inspects? (\textbf{\sk{अवेक्षते}}, \emph{avekṣate}.) Devises? (\textbf{\sk{उद्भावयति}}, \emph{udbhāvayati}.) Studying? (\textbf{\sk{अध्ययनम्}}, \emph{adhyayanam}.) A topic? (\textbf{\sk{विषयः}}, \emph{viṣayaḥ}.) A branch of learning? (\textbf{\sk{शास्त्रम्}}, \emph{śāstram}.) Studies? (\textbf{\sk{अधीते}}, \emph{adhīte}.) Practises? (\textbf{\sk{अभ्यस्यति}}, \emph{abhyasyati}.)
\end{tcolorbox}
